% =====================================================================
% Слайд 15 — Симулятор солиста и рендера
% =====================================================================
\begin{frame}{Симулятор солиста и рендера}
  \footnotesize

  Для PPO нужны эпизоды --- (state, action, reward) в больших
  количествах. На реальных ASAP-записях их мало ($\sim$1000
  исполнений). Решение --- симулятор.

  \begin{columns}[T,onlytextwidth]
    \column{0.50\textwidth}
    \textbf{Симулятор солиста:}
    \begin{itemize}\setlength{\itemsep}{0.2em}
      \item база --- реальное ASAP-исполнение или партитура;
      \item параметрическое rubato (Repp 1995);
      \item сэмплирование ошибок (Nakamura 2015): вставка / замена /
            пропуск;
      \item диверсивность $\Rightarrow$ агент не переобучается на
            конкретных записях.
    \end{itemize}

    \vspace{0.4em}
    \textbf{Симулятор рендера:}
    \begin{itemize}\setlength{\itemsep}{0.2em}
      \item детерминированное проигрывание orchestra MIDI с темпом
            $a_t$;
      \item $t^{\text{render}}_t$ --- время события после temp-rescaling.
    \end{itemize}

    \column{0.50\textwidth}
    \centering
    \resizebox{\linewidth}{!}{%
      % СХЕМА 11 — Замкнутый цикл симулятора
\begin{tikzpicture}[
  every node/.style={font=\scriptsize, align=center},
  box/.style={draw, thick, minimum height=0.85cm, minimum width=1.9cm,
              fill=white, rounded corners=1.5pt, inner sep=3pt},
  rlbox/.style={box, line width=1pt, fill=black!4},
  arr/.style={-{Latex[length=2mm]}, thick},
]
  \node[box]   (sim)   at (0, 1.5)    {Симулятор солиста\\(ASAP + rubato + errors)};
  \node[box]   (tr)    at (4.5, 1.5)  {HybridScore\\Follower};
  \node[rlbox] (pi)    at (4.5, -0.3) {$\pi_\theta$};
  \node[box]   (rend)  at (0, -0.3)   {Симулятор рендера\\(orchestra MIDI $\times a_t$)};
  \node[rlbox] (rew)   at (-2.7, 0.6) {Reward $r_t$\\\scriptsize(формула 1)};

  \draw[arr] (sim) -- node[above] {$(pitch, t^{\text{perf}})$} (tr);
  \draw[arr] (tr)  -- node[right] {$s_t$} (pi);
  \draw[arr] (pi)  -- node[below] {$a_t$} (rend);
  \draw[arr] (rend.west) -| node[below, near start] {$t^{\text{render}}$} (rew);
  \draw[arr] (rew.north) |- (pi.west);
\end{tikzpicture}
%
    }\\[0.4em]
    {\scriptsize\itshape\color{black!75}
     Замкнутый цикл: симулятор $\to$ трекер $\to$\\
     политика $\to$ рендер $\to$ reward $\to$ обратно.}
  \end{columns}

  \vspace{0.3em}
  {\centering\footnotesize
   $t^{\text{render}}_t - t^{\text{perf}}_t$ доступен \rlhi{без}
   аудио-записи оркестра --- достаточно solo MIDI $+$ orchestra MIDI.\par}
\end{frame}
